\subsection{半单代数的对偶}

设 \(\Theta^{\mathrm{ss}}(\mathrm{A})\) 是左 A-模 \(\Theta (\mathrm{A})\) 的底座，即（VIII, 第 65 页）\(\Theta (\mathrm{A})\) 的最大半单子模。我们用 \(\mathcal{S}_{\mathrm{K}}\) 表示在 K 上有限维的单（左）A-模的类所成的集。当 A 是 K 上一个有限次数的半单代数时，我们有 \(\mathrm{A}^{*} = \Theta (\mathrm{A}) = \Theta^{\mathrm{ss}}(\mathrm{A})\)，因为每个左 A-模都是半单的（VIII, 第 138 页，命题 4）。

定理 1. —— a) 集 \(\Theta^{\mathrm{ss}}(\mathrm{A})\) 由 A 的有限维半单表示的系数组成。它是 \(\mathrm{A}^*\) 的一个 (A,A)-子双模。

\begin{enumerate}
\def\labelenumi{\alph{enumi})}
\setcounter{enumi}{1}
\item
  对每个 \(S \in \mathcal{S}_{\mathrm{K}}\)，\(\Theta(A)\) 的 S 型同型分量等于 \(\Theta_{\mathrm{S}}(\mathrm{A})\)。左 A-模 \(\Theta^{\mathrm{ss}}(\mathrm{A})\) 是诸子模 \(\Theta_{\mathrm{S}}(\mathrm{A})\) 的直和，其中 S 取遍 \(\mathcal{S}_{\mathrm{K}}\)。
\item
  对 \(\mathcal{S}_{\mathrm{K}}\) 中每个 S，右 A-模 \(S^{*}\) 是单的，且 \(\Theta_{\mathrm{S}}(\mathrm{A})\) 是 \(S^{*}\) 型的同型分量（右 A-模 \(\Theta(\mathrm{A})\) 的）。把 S 映为 \(\operatorname{cl}(\mathrm{S}^{*})\) 的映射是从 \(\mathcal{S}_{\mathrm{K}}\) 到 K 上有限维单 \(A^{o}\) -模的类集的一个双射。
\item
  视为右 A-模时，\(\Theta (\mathrm{A})\) 的底座是 \(\Theta^{\mathrm{ss}}(\mathrm{A})\)。
\end{enumerate}

设 E 是一个在 K 上有限维的半单 A-模。E 的每个系数属于从 E 到 \(A^{*}\) 的某个 A-线性映射的像（VIII, 第 378 页，命题 2），因而属于 \(\Theta^{\mathrm{ss}}(\mathrm{A})\)。反之，设 \(f \in \Theta^{\mathrm{ss}}(\mathrm{A})\)。则 A-模 Af 在 K 上有限维且是半单的。由 VIII, 第 367 页的注，f 是 Af 的一个系数。

对 \(\mathcal{S}_{\mathrm{K}}\) 中每个 S，\(\Theta(\mathrm{A})\) 的 S 型同型分量由从 S 到 \(A^{*}\) 的 A-线性映射的像生成；因而由 VIII, 第 378 页，命题 2, a)，它等于 \(\Theta_{\mathrm{S}}(\mathrm{A})\)。另一方面，若 S 是一个在 K 上非有限维的单 A-模，则 \(\Theta(\mathrm{A})\) 的 S 型同型分量为零：事实上，\(\Theta(\mathrm{A})\) 的每个单子模都是单生成的，因此在 K 上有限维。由于 \(\Theta(\mathrm{A})\) 的底座是它的诸同型分量的直和（VIII, 第 65 页，命题 4），这就证明了 b)。

设 S 是一个在 K 上有限维的单 A-模。由于 K-向量空间 S 不约化为 0，\(S^{*}\) 亦然。设 E 是右 A-模 \(S^{*}\) 的一个子模；它在 S 中的正交 \(E'\) 是 S 的一个 A-子模。由于 S 是单的，我们或者有 \(E' = 0\)（此情形 \(E = S^{*}\)），或者有 \(E' = S\)（此情形 E = 0）。因此 \(S^{*}\) 是一个单右 A-模。

我们有 \(\Theta (\mathrm{A}^{\circ}) = \Theta (\mathrm{A})\)（VIII, 第 376 页）；我们把右 A-模等同于左 \(\mathrm{A}^{\circ}\) -模。由于 K 上每个有限维向量空间同构于它的双对偶，上述论证证明映射 \(\mathrm{S}\mapsto \mathrm{cl}(\mathrm{S}^{*})\) 是从 \(\mathcal{S}_{\mathrm{K}}\) 到 K 上有限维单 \(\mathrm{A}^{\circ}\) -模的类集的一个双射。现在，对 \(\mathcal{S}_{\mathrm{K}}\) 中的 S，\(\Theta (\mathrm{A}^{\circ})\) 的 \(\mathrm{S}^*\) 型同型分量等于 \(\Theta_{\mathrm{S}^*}(\mathrm{A}^{\circ})\)（由应用于代数 \(\mathrm{A}^{\circ}\) 的断言 b)），且我们有 \(\Theta_{\mathrm{S}^*}(\mathrm{A}^{\circ}) = \Theta_{\mathrm{S}}(\mathrm{A})\)。断言 c) 与 d) 随即得出。

推论 1. —— K 上每个有限维单 A-模都同构于 \(\Theta(A)\) 的一个 A-子模。

这由定理 1, b) 得出，因为对每个非零 A-模 E 我们有 \(\Theta_{\mathrm{E}}(\mathrm{A}) \neq 0\)。

推论 2. —— 若 K 上两个有限维单 A-模有一个公共的非零系数，则它们同构。

设 \(S_{1}\) 与 \(S_{2}\) 是 K 上有公共非零系数的有限维单 A-模。则我们有 \(\Theta_{S_{1}}(A) \cap \Theta_{S_{2}}(A) \neq 0\)。而 \(\Theta_{S_{1}}(A)\) 是 \(S_{1}\) 型的同型模，\(\Theta_{S_{2}}(A)\) 是 \(S_{2}\) 型的同型模。故 \(S_{1}\) 同构于 \(S_{2}\)。

由定理 1，\(\Theta^{\mathrm{ss}}(\mathrm{A})\) 是 \(\mathrm{A}^*\) 的一个 (A,A)-子双模，即诸 (A,A)-双模 \(\Theta_{\mathrm{S}}(\mathrm{A})\)（S 取遍 \(\mathcal{S}_{\mathrm{K}}\)）的直和。这些双模两两不同构，因为它们作为左 A-模已经不同构。

固定 \(\mathcal{S}_{\mathrm{K}}\) 中的一个 S。用 D 表示 \(\operatorname{End}_{\mathrm{A}}(\mathrm{S})\) 的反体，并把 S 视为一个右 D-模，\(S^{*}\) 视为一个左 D-模。在这些约定下，S 是一个 (A, D)-双模，\(S^{*}\) 是一个 (D, A)-双模，\(S \otimes_{D} S^{*}\) 是一个 (A, A)-双模。

命题 4. —— a) 存在一个群同态 \(\lambda_{\mathrm{S}}\)，它从 \(\mathrm{S} \otimes_{\mathrm{D}} \mathrm{S}^{*}\) 到 \(\Theta_{\mathrm{S}}(\mathrm{A})\)，由

\[
\lambda_ {\mathrm{S}} (x \otimes x ^ {*}) = c _ {\mathrm{S}} (x, x ^ {*})\tag{6}
\]

（对 \(x \in S\) 与 \(x^{*} \in S^{*}\)）刻画。这个映射是 (A, A)-双模的一个同构。

\begin{enumerate}
\def\labelenumi{\alph{enumi})}
\setcounter{enumi}{1}
\tightlist
\item
  (A, A)-双模 \(\Theta_{\mathrm{S}}(\mathrm{A})\) 是单的。
\end{enumerate}

映射 \(\gamma_{\mathrm{S}}: \mathrm{S} \otimes_{\mathrm{K}} \mathrm{S}^{*} \to \Theta_{\mathrm{S}}(\mathrm{A})\)（由公式 \(\gamma_{\mathrm{S}}(x \otimes x^{*}) = c_{\mathrm{S}}(x, x^{*})\) 刻画）是 (A, A)-线性的且是满射的，并满足 \(\gamma_{\mathrm{S}}(xd \otimes x^{*}) = \gamma_{\mathrm{S}}(x \otimes dx^{*})\)（对 \(x \in \mathrm{S}\)、\(x^{*} \in \mathrm{S}^{*}\) 与 \(d \in \mathrm{D}\)）。因此它定义了一个由公式 (6) 刻画的满射 (A, A)-线性映射 \(\lambda_{\mathrm{S}}: \mathrm{S} \otimes_{\mathrm{D}} \mathrm{S}^{*} \to \Theta_{\mathrm{S}}(\mathrm{A})\)。

我们来证明 \(S \otimes_{D} S^{*}\) 是一个单 (A, A)-双模。由 VIII, 第 63 页，推论 2，\(S \otimes_{D} S^{*}\) 的每个 (A, A)-子双模都具有形式 \(S \otimes_{D} H\)，其中 H 是 \(S^{*}\) 的一个 (D, A)-子双模。由于 \(S^{*}\) 是一个单右 A-模（VIII, 第 380 页，定理 1, c)），我们有 H = 0 或 \(H = S^{*}\)；断言得证。

同态 \(\lambda_{S}\) 是 (A, A)-线性的且非零；由于 \(S \otimes_{D} S^{*}\) 是一个单双模，\(\lambda_{S}\) 是单射的（VIII, 第 47 页，命题 2, a)）。

注. —— 当体 K 是代数闭的时，由 VIII, 第 47 页，定理 1，我们有 D = K，且 \(\lambda_{S}\) 是从 \(S \otimes_{K} S^{*}\) 到 \(\Theta_{\mathrm{S}}(\mathrm{A})\) 的 (A, A)-双模的一个同构。
% semantic-fragments: legacy-input-seam
\relax{}
